\section{Training}
\label{sec:training}

The training process underwent several iterative refinements regarding dataset curation, hardware constraints, class imbalance mitigation, and regularization strategies to prevent overfitting.

\subsection{Dataset Curation and Challenges}
Initial experiments utilizing a small subset of data revealed a significant performance gap compared to fully supervised detectors, primarily because our initial subset of 20 images per breed was insufficient. Although the Cat-API provided access to a large volume of entries, the data suffered from inconsistent labeling, yielding only about 12 usable images per breed. Manually augmenting or labeling an additional eight images per class proved unfeasible. 

To resolve this, we integrated the Oxford-IIIT Pet Dataset \cite{parkhi2012catsdogs}, which expanded our image pool by 2,744 relevant samples. During inspection, we noted that \textit{Tabby} and \textit{Tiger\_cat} shared high visual similarity, with \textit{Tabby} representing a fur pattern rather than a distinct biological breed; consequently, we standardized our target classes around \textit{Tiger\_cat}. Furthermore, major dog breeds such as the \textit{Golden\_Retriever}, \textit{German\_Shepherd}, \textit{Siberian\_Husky}, \textit{Dalmatian}, and \textit{Rottweiler} were absent from the Oxford dataset. To close this gap, we incorporated the Stanford Dogs Dataset \cite{khosla2011novel}, providing 150 images per specified breed to finalize our primary training repository.

To address persistent ambiguities with tabby and tiger cats, we supplemented these specific classes with 30 additional images sourced from Unsplash and Pixabay. Because initial classification performance remained unsatisfactory for these categories, we increased the density to 50 images per breed, successfully pushing overall performance above our target threshold of 80\% accuracy.

During subsequent training runs, performance occasionally degraded due to dataset shifts. A structural audit revealed that Dalmatian samples had been temporarily excluded. Since earlier API endpoints lacked Dalmatian data, we retrieved an independent archive of 359 Dalmatian images and integrated it into the dataset. This addition yielded a substantial performance jump, raising accuracy from 17\% to 69\% after 10 epochs. Further attempts to generate synthetic images via AI or web-scrape additional uncurated data failed due to excessive semantic noise (e.g., queries for the ``Bombay'' cat breed predominantly returned bottles of Bombay Sapphire Gin) and high manual cleaning overhead.

\subsection{Hardware Optimization and Batch Dynamics}
Initial training configurations suffered from hardware limitations on our shared GPU infrastructure (7.6 GiB VRAM). Memory consumption on the GPU does not scale linearly with image resolution; rather, it grows quadratically because images possess two spatial dimensions (width $\times$ height). For instance, an input resolution of $224 \times 224$ contains 50,176 pixels, whereas an increased resolution of $384 \times 384$ scales this to 147,456 pixels. 

Increasing the image size from 224 to 384 forced the GPU to store nearly triple the amount of pixel data and intermediate activations per sample. In deeper architectures such as EfficientNetV2 \cite{tan2021efficientnetv2}, this memory footprint multiplies rapidly across layers. Combined with a batch size of 32, the network attempted to allocate simultaneous activation tensors for $32 \times 147,456 \text{ pixels} \times \text{layers}$ during the forward pass, triggering frequent Out-Of-Memory (OOM) crashes. Halving the batch size to 16 and reducing the input resolution to $224 \times 224$ successfully stabilized VRAM allocation.

\subsection{Handling Class Imbalance via Loss Weighting}
To counteract severe class imbalances where rare breeds were dominated by frequent classes, we implemented a dynamic class-weighting mechanism inside PyTorch's cross-entropy loss function. Without intervention, models tend to ignore minority classes to artificially lower the global loss.

We computed individual class weights $w_i$ for each class $i$ using the standard inverse frequency formulation:
\begin{equation}
    w_i = \frac{N_{\text{total}}}{C \cdot n_i}
\end{equation}
where $N_{\text{total}}$ represents the total number of training samples, $C$ denotes the total number of classes (10 per species split), and $n_i$ is the number of images available in class $i$. Frequent classes received lower weights (e.g., $\approx 0.5$), while underrepresented classes received higher weights (e.g., $\approx 5.0$). 

During backpropagation, misclassifications of minority-class samples are penalized proportionally higher ($Loss_{\text{weighted}} = Loss_{\text{standard}} \cdot w_i$), forcing the network gradients to correct errors on rare breeds aggressively.

\subsection{Optimization Dynamics and Overfitting Control}
Training deep neural networks requires careful tuning of learning rates and convergence monitoring. Extremely low learning rates trap models in local minima or saddle points, where flat loss landscapes yield vanishing gradients that starve the optimization process of necessary momentum. Conversely, unmanaged training durations lead directly to overfitting, where models memorize training samples rather than generalizing to validation splits.

To automate optimal checkpoint selection, we implemented model checkpointing combined with validation-based early stopping (with a patience threshold). Rather than selecting the final epoch's weights, the pipeline saves model weights strictly when validation accuracy improves. 

Furthermore, to prevent training-validation divergence, we enforced strict regularization practices, combining large dataset volumes, multi-epoch training with decaying learning rates (Cosine Annealing), high dropout layers ($p = 0.5$), and custom spatial augmentations such as \texttt{CropMix}.

\bibliographystyle{plain}
\bibliography{main}

\end{document}